\documentclass{tufte-handout}

\usepackage{style}

\begin{document}

\tma{04} % TMA number

\begin{question}

\begin{equation*}
f_{n}(x) = 
\begin{cases}
\sin(\pi x), \text{ if } 2n \leq x < (2n+1),\\
0, \text{ otherwise. }
\end{cases}
\end{equation*}

\qpart

\includegraphics[scale=0.5]{name}

\vspace{2cm}

\includegraphics[scale=0.5]{name}

\vspace{2cm}

\includegraphics[scale=0.5]{name}

\vspace{2cm}

\qpart

To test of a pointwise convergence of \( f_n(x) \) to \( f(x) \), we need to consider different cases for \( x \).

\begin{align*}
\text{Case 1: } & x \in [0,1) \\
& f_n(x) = \sin(\pi x) \text{ for all } n \geq 0 \\
& \lim\limits_{n \to \infty} f_n(x) = \sin(\pi x) \\
& \therefore f(x) = \sin(\pi x) \\[12pt]
\text{Case 2: } & x \in [1,2) \\
& f_n(x) = 0 \text{ for all } n \geq 0 \\
& \lim\limits_{n \to \infty} f_n(x) = 0 \\
& \therefore f(x) = 0 \\[12pt]
\text{Case 3: } & x \in [2,3) \\
& f_n(x) = \sin(\pi x) \text{ for all } n \geq 1 \\
& \lim\limits_{n \to \infty} f_n(x) = \sin(\pi x) \\
& \therefore f(x) = \sin(\pi x) \\[12pt]
\text{Case 4: } & x \in [3,4) \\
& f_n(x) = 0 \text{ for all } n \geq 1 \\
& \lim\limits_{n \to \infty} f_n(x) = 0 \\
& \therefore f(x) = 0 \\[12pt]
& \vdots \\[12pt]
\text{Case k: } & x \in [k, k+1) \\
& f_n(x) = 
\begin{cases}
\sin(\pi x), \text{ if } k \text{ is even and } n \geq k/2,\\
0, \text{ if } k \text{ is odd and } n \geq (k-1)/2.
\end{cases} \\
& \lim\limits_{n \to \infty} f_n(x) = 
\begin{cases}
\sin(\pi x), \text{ if } k \text{ is even},\\
0, \text{ if } k \text{ is odd. }
\end{cases} \\
& \therefore f(x) = 
\begin{cases}
\sin(\pi x), \text{ if } k \text{ is even},\\
0, \text{ if } k \text{ is odd. }
\end{cases}     
\end{align*}

Thus we have that the pointwise limit function \( f(x) \) is given by;
\begin{equation*}
f(x) = 
\begin{cases}
\sin(\pi x), \text{ if } \lfloor x \rfloor \text{ is even},\\
0, \text{ if } \lfloor x \rfloor \text{ is odd. }
\end{cases}
\end{equation*} 



\vspace{3cm}

\qpart



\end{question}

\end{document}